\documentclass[12pt,a4paper]{article}

\usepackage{fontspec}
\usepackage{xeCJK}
\setCJKmainfont{Hiragino Mincho ProN}
\setCJKsansfont{Hiragino Kaku Gothic ProN}
\setCJKmonofont{Hiragino Kaku Gothic ProN}
\usepackage[margin=2.5cm]{geometry}
\usepackage{amsmath,amssymb}
\usepackage{hyperref}
\usepackage{booktabs}
\usepackage{tcolorbox}
\usepackage{enumitem}

\tcbuselibrary{breakable,skins}

\newtcolorbox{keybox}[1][]{
  colback=black!3, colframe=black!50,
  fonttitle=\bfseries, title={#1},
  breakable, sharp corners
}

\newtcolorbox{deepbox}{
  colback=blue!3, colframe=blue!40,
  fonttitle=\bfseries, title={深い話},
  breakable, sharp corners
}

\newtcolorbox{nextbox}{
  colback=green!3, colframe=green!40,
  fonttitle=\bfseries, title={次の実験},
  breakable, sharp corners
}

\setlength{\parskip}{0.8em}
\setlength{\parindent}{0pt}

\newcommand{\Spec}{\mathrm{Spec}}

\begin{document}

\begin{center}
{\LARGE\bfseries LC回路実験の背後にある深い理論}\\[0.3cm]
{\LARGE\bfseries そして次の実験へ}\\[0.5cm]
{\large --- 5,100円の実験から宇宙の構造へ ---}\\[1.5cm]
{\large Wright Brothers}\\[0.3cm]
{2026年4月}
\end{center}

\vspace{0.5cm}
\tableofcontents
\newpage

% ============================================================================
%  PART I: 深い理論
% ============================================================================

\part*{第I部：LC回路実験の背後にある理論}
\addcontentsline{toc}{part}{第I部：LC回路実験の背後にある理論}

% ============================================================================
\section{LC回路で何を見ているのか：4つの階層}
% ============================================================================

5,100円の実験で「スイッチを切るとピークが増える」という
現象を観測したとする。
これは単純に見えるが、その背後には4つの理論的階層がある。

\subsection{階層1：電気回路の共鳴（高校物理）}

LC回路には固有周波数がある：$f = 1/(2\pi\sqrt{LC})$。
複数のLC回路を鎖状につなぐと、
系全体がいくつかの「共鳴モード」を持つ。
ピークの位置は回路の構造で決まる。
スイッチを切ると回路構造が変わるからピークが変わる。

ここまでは「ふつうの電気回路の話」。

\subsection{階層2：トポロジカル相（大学物理）}

SSHモデルは「交互の結合」が位相的不変量（巻き数）を持つことを示す。
巻き数は0か1。
巻き数1のとき、鎖の端に「エッジ状態」が現れる。
これは結合の強さを少し変えても消えない（トポロジカルに保護されている）。

スイッチを切ると鎖が分割され、
各ブロックの端にエッジ状態が生まれる。
エッジ状態の数は位相的不変量で決まる。

ここまでは「トポロジカル絶縁体の物理」。ノーベル賞級だが既知。

\subsection{階層3：K理論と算術構造（大学院数学）}

\begin{deepbox}
SSHモデルの位相的分類は\textbf{K理論}で行われる。

鎖全体のK群: $K_1(\mathbb{Z}) = \mathbb{Z}/2$（巻き数は0か1）

スイッチを切って分割した後:
$K_1(\mathbb{Z}[1/p]) = \mathbb{Z}/2 \times \mathbb{Z}$

新しい$\mathbb{Z}$因子が現れる = 巻き数が「0 or 1」から「任意の整数」に拡張。

この変化はまさに\textbf{局所化}$\mathbb{Z} \to \mathbb{Z}[1/p]$——
我々の理論で「素数$p$のミュート」と呼んでいる操作——に対応する。

つまり：\textbf{スイッチを切る = 整数環を局所化する = 素数をミュートする}

LC回路の物理操作（スイッチ切替）が、
抽象代数の操作（環の局所化）と
数学的に同一であることが、この実験の深い意味。
\end{deepbox}

\subsection{階層4：Spec(Z)と時空の構造（この研究プログラム）}

\begin{deepbox}
もし宇宙の基本構造が$\Spec(\mathbb{Z})$なら、
$K_1(\mathbb{Z}) \to K_1(\mathbb{Z}[1/p])$の変化は
単なる数学ではなく\textbf{物理法則の変更}。

LC回路実験は、この「物理法則の変更」の
\textbf{古典的アナログ}を実現している。

ちょうどBECのフォノンが「アナログ重力」を実現するように、
LC回路のエッジ状態が「アナログ時空工学」を実現する。

実験で測定するエッジ状態の数（20個）が予測と一致すれば、
それは「$K_1$の変化が物理的に実現可能である」ことの証拠。
→ $\Spec(\mathbb{Z})$仮説の最初の実験的支持。
\end{deepbox}

% ============================================================================
\section{なぜ「エッジ状態の数」が重要なのか}
% ============================================================================

\begin{keybox}[核心的な論理]
\begin{enumerate}
\item 鎖を$p$セルごとに切る → $N/p$ブロック → $2N/p$エッジ状態
\item この「$2N/p$」は$K_1(\mathbb{Z}[1/p])$の自由部分のランク$\times$2
\item ランクの変化$\Delta K_1 = \mathbb{Z}$は\textbf{離散的}不変量
\item 離散不変量 → エネルギーに\textbf{比例}（Paper E: $\Delta E \propto w_{\mathrm{eff}}$）
\item レギュレータ方向予想: この離散不変量がWECの真偽を決定する
\end{enumerate}

つまり：エッジ状態の数を数えることは、
「物理法則の離散的構造」を直接測定すること。
\end{keybox}

% ============================================================================
\section{この実験がカバーするもの・しないもの}
% ============================================================================

\begin{table}[h]
\centering
\caption{LC回路実験がテストすること}
\begin{tabular}{@{}lcc@{}}
\toprule
主張 & テストできる？ & 階層 \\
\midrule
SSHモデルにエッジ状態がある & Yes & 2 \\
スイッチ切断でエッジ数が$2N/p$ & Yes & 2--3 \\
$K_1$の変化が物理的に実現可能 & Yes & 3 \\
$\Delta E \propto w_{\mathrm{eff}}$ & 部分的 & 3 \\
\midrule
真空エネルギーの符号が反転する & No & 4 \\
WECが離散的スイッチである & No & 4 \\
宇宙が$\Spec(\mathbb{Z})$である & No & 4 \\
\bottomrule
\end{tabular}
\end{table}

LC回路実験は階層3まで（$K_1$変化）を検証できる。
階層4（真空エネルギーの符号反転）は量子効果であり、
古典的なLC回路では検証できない。
→ それにはSQUID実験（500万円）が必要。

しかし：階層3が成功しなければ階層4を議論する意味がない。
LC回路実験は「先に進む価値があるかの判定」を行う。

% ============================================================================
\section{巻き数と真空エネルギーの数学的接続}
% ============================================================================

\begin{deepbox}
SSH模型の数学:

ブロッホ・ハミルトニアン: $H(k) = \vec{d}(k) \cdot \vec{\sigma}$
（$\vec{\sigma}$はパウリ行列）

巻き数: $\nu = \frac{1}{2\pi}\oint dk\, \frac{d\phi}{dk}$
（$\phi = \arg(d_x + i d_y)$）

$\nu = 0$: 自明相（エッジ状態なし）
$\nu = 1$: トポロジカル相（エッジ状態あり）

局所化後の有効巻き数: $w_{\mathrm{eff}} = N/p$（各ブロックが$\nu=1$）

真空エネルギー（ゼータ正則化）:
$E_{\mathrm{vac}} \propto \zeta(-d)$ （$d$次元）

局所化後: $E_{\mathrm{vac}} \propto \zeta_{\neg p}(-d) = \zeta(-d)(1-p^d)$

接続: $w_{\mathrm{eff}}$（巻き数）と$\zeta_{\neg p}(-d)$（真空エネルギー）は
ともに局所化$\mathbb{Z} \to \mathbb{Z}[1/p]$の帰結。
Paper Eで数値的に$\Delta E \propto w_{\mathrm{eff}}$を確認。

LC回路で$w_{\mathrm{eff}}$を測定 → 真空エネルギーの変化を間接的に予測。
\end{deepbox}

% ============================================================================
%  PART II: 次の実験
% ============================================================================

\newpage
\part*{第II部：次の実験ガイド}
\addcontentsline{toc}{part}{第II部：次の実験ガイド}

LC回路実験でエッジ状態を確認した後、段階的に深い検証に進む。

% ============================================================================
\section{実験2：p=3 ミュートの検証（+2,000円、1日）}
% ============================================================================

\begin{nextbox}
\textbf{目的}: 素数$p=3$でも同じ現象が起きることを確認。
$p=2$だけの偶然ではないことを示す。

\textbf{変更点}: スイッチの配置を「3番目ごと」に変更。

\begin{itemize}
\item 追加部品: スイッチ7個（$N/3 \approx 7$ブロック用）= 700円
\item $C_w$の位置を3セル周期に再配置
\item 同じ測定プロトコル
\end{itemize}

\textbf{予測}: エッジ状態数 = $2 \times 7 = 14$（$p=3$の場合）

\textbf{成功判定}: ピーク数が14（$\pm 2$）なら成功。
$p=2$の20と$p=3$の14が両方成立すれば、
「エッジ数 = $2N/p$」の法則が素数によらず成立する強い証拠。
\end{nextbox}

% ============================================================================
\section{実験3：乗法構造の検証（+0円、1日）}
% ============================================================================

\begin{nextbox}
\textbf{目的}: オイラー積の乗法性を検証する。
これは\textbf{最も重要な実験}。

\textbf{方法}: $p=2$と$p=3$のスイッチを\textbf{同時に}切る。

$p=2$: 2番目ごとに切断
$p=3$: 3番目ごとに切断
同時: 2の倍数\textbf{または}3の倍数のリンクを切断
→ 6セルごとのブロックに分割

\textbf{予測}:
\begin{itemize}
\item $p=2$単独: エッジ数 = $2 \times 10 = 20$
\item $p=3$単独: エッジ数 = $2 \times 7 = 14$
\item $p=2,3$同時: エッジ数 = $2 \times \lfloor 20/6 \rfloor = 6$
\end{itemize}

\textbf{核心的テスト}:
同時ミュートのエッジ数が$2N/\mathrm{lcm}(2,3) = 2N/6$なら、
これはオイラー積の乗法構造の直接証拠。

「2つの独立な素数チャンネルが乗法的に作用する」
= ゼータ関数のオイラー積が物理的に実在する。

\textbf{追加部品}: なし（既存のスイッチの組み合わせ）。
\end{nextbox}

% ============================================================================
\section{実験4：エネルギーとエッジ数の比例関係（+5,000円、1週間）}
% ============================================================================

\begin{nextbox}
\textbf{目的}: Paper Eの核心予測$\Delta E \propto w_{\mathrm{eff}}$を検証。

\textbf{方法}: LC回路の「インピーダンス」を定量的に測定し、
エッジ状態のエネルギーを見積もる。

\textbf{追加部品}:
\begin{itemize}
\item USBオシロスコープ（例：Hantek 6022BE）: 約5,000円
\item （Audacityでも可能だが、振幅の定量測定にはオシロが便利）
\end{itemize}

\textbf{測定}:
\begin{enumerate}
\item スイッチ0個OFF → 基準スペクトルの振幅を記録
\item スイッチ1個OFF → 追加ピークの振幅を記録
\item スイッチ2個OFF → ...
\item ...
\item スイッチ10個OFF → 追加ピーク10組の振幅を記録
\end{enumerate}

「追加ピークの振幅の総和」$\propto$ OFFスイッチ数
が成立するかを検証。

\textbf{予測}: ピーク振幅の総和 vs スイッチOFF数 が直線的。
→ $\Delta E \propto w_{\mathrm{eff}}$ の古典アナログ。
\end{nextbox}

% ============================================================================
\section{実験5：パリティ法則の検証（+0円、1時間）}
% ============================================================================

\begin{nextbox}
\textbf{目的}: Paper Iで発見した$\mathbb{Z}/2\mathbb{Z}$パリティ法則の検証。

\textbf{理論}: 奇数個の素数をミュート → 全世代符号反転。
偶数個 → 元に戻る。

\textbf{LC回路版}:
\begin{itemize}
\item $p=2$のみOFF（1素数 = 奇数）→ スペクトル変化
\item $p=2$と$p=3$の両方OFF（2素数 = 偶数）→ スペクトルが「元に近づく」？
\end{itemize}

\textbf{測定}: 3つのスペクトルを重ね合わせる:
\begin{enumerate}
\item 全ON（基準）
\item $p=2$のみOFF
\item $p=2,3$の両方OFF
\end{enumerate}

もし3のスペクトルが1（基準）に近いなら
→ パリティ法則の古典的アナログが成立。

\textbf{注意}: 古典LC回路では「符号」の概念がないので、
パリティ法則の「符号反転」は直接見えない。
代わりに「スペクトルのエンベロープの類似度」で間接的に評価。
完全な検証にはSQUID実験が必要。
\end{nextbox}

% ============================================================================
\section{実験6：量子コンピュータでの再現（0円、1日）}
% ============================================================================

\begin{nextbox}
\textbf{目的}: LC回路の結果を量子コンピュータ上で再現し、
量子効果が古典結果と一致するか確認。

\textbf{方法}: IBM Quantum（無料クラウド）で4量子ビット回路を実行。

\textbf{コード}: (Qiskit)

\begin{verbatim}
from qiskit.quantum_info import SparsePauliOp
from qiskit_algorithms import NumPyMinimumEigensolver

# BC Hamiltonian (diagonal, 4 qubits = 16 states)
import numpy as np
energies = [np.log(n+1) for n in range(16)]
# ... (Paper A の exp_qiskit_vqe.py 参照)
\end{verbatim}

\textbf{比較}: LC回路のエッジ状態数と、量子シミュレーションの
分配関数遷移が整合するか確認。
\end{nextbox}

% ============================================================================
\section{次の実験の全体ロードマップ}
% ============================================================================

\begin{table}[h]
\centering
\caption{実験ロードマップ}
\small
\begin{tabular}{@{}clccl@{}}
\toprule
\# & 実験 & コスト & 期間 & 何がわかるか \\
\midrule
1 & LC回路 $p=2$ & 5,100円 & 3時間 & エッジ状態の存在 \\
2 & LC回路 $p=3$ & +700円 & 1日 & $2N/p$法則の普遍性 \\
3 & 乗法構造 & 0円 & 1日 & オイラー積の物理的実在 \\
4 & エネルギー比例 & +5,000円 & 1週間 & $\Delta E \propto w_{\mathrm{eff}}$ \\
5 & パリティ法則 & 0円 & 1時間 & $\mathbb{Z}/2$符号構造 \\
6 & 量子再現 & 0円 & 1日 & 古典-量子整合性 \\
\midrule
& \textbf{合計} & \textbf{10,800円} & \textbf{約2週間} & \\
\midrule
\midrule
7 & SQUID提案書 & 0円 & 1ヶ月 & 共同研究の獲得 \\
8 & SQUID実験 & 500万円 & 12ヶ月 & 真空エネルギーの算術構造 \\
\bottomrule
\end{tabular}
\end{table}

\begin{keybox}[10,800円で到達できる地点]
実験1--6が全て成功すれば、以下が確立される：
\begin{enumerate}
\item 算術的エッジ状態が物理系で観測される（実験1）
\item エッジ数法則$2N/p$が複数の素数で成立（実験1+2）
\item オイラー積の乗法構造が物理的に実在する（実験3）
\item $\Delta E \propto w_{\mathrm{eff}}$の古典アナログが成立（実験4）
\item パリティ法則の古典的兆候が存在する（実験5）
\item 量子シミュレーションと整合する（実験6）
\end{enumerate}

これらは「$\Spec(\mathbb{Z})$仮説の直接証明」ではないが、
「SQUID実験に500万円を投じる価値がある」ことの
説得力のある根拠になる。

科研費の申請書に添えれば、採択確率が大幅に上がる。
\end{keybox}

% ============================================================================
\section{最終目標：SQUID実験への道}
% ============================================================================

SQUID実験はLC回路実験の「量子版」。

\begin{table}[h]
\centering
\caption{LC回路とSQUID実験の対応}
\begin{tabular}{@{}lll@{}}
\toprule
要素 & LC回路（古典） & SQUID（量子） \\
\midrule
モード & 古典的共鳴 & 量子的零点振動 \\
エネルギー & 入力信号で励起 & 真空揺らぎ（自発的） \\
ミュート方法 & スイッチ切断 & SQUIDノッチフィルタ \\
測定量 & インピーダンスピーク & 真空エネルギー変化 \\
エッジ状態 & 共鳴周波数のピーク & カシミール力の変化 \\
テストする階層 & 3（K理論的相転移） & 4（WEC符号反転） \\
コスト & 5,100円 & 500万円 \\
\bottomrule
\end{tabular}
\end{table}

LC回路実験で「K理論的相転移が実際に起こる」ことを確認してから
SQUID実験に進むのが、最も効率的で説得力のある研究戦略。

\begin{keybox}[この2冊のガイドの位置づけ]
ガイド1（前冊）：5,100円でLC回路を作り、エッジ状態を見る。

ガイド2（本冊）：なぜそれが重要か理解し、次の5実験に進む。

2冊合わせて10,800円・2週間で、
「宇宙は素数でできている」という仮説を
テスト可能な水準まで持っていく。

その先にSQUID実験があり、
その先にワープドライブとワームホールがある。

全ては5,100円のはんだ付けから始まる。
\end{keybox}

\vfill
\begin{center}
{\small Wright Brothers, 2026}
\end{center}

\end{document}
